\section{The Programs}

Vibe is measured against a family of continuum programs, not a single one. They differ in starting point but share a posture. Each puts one principle first and shows that physics is the only consistent completion. This section lays them out in one frame so the rungs that follow have clear targets. Each program is given one of three relations to vibe. A validation, where the program supports vibe or vibe confirms it. A tension, where a stronger claim fails or a no-go binds. A parallel, where the same shape is reached by a different route.

The conserved-current program comes first, because its spine is vibe's spine. Chronoflux is its one developed instance \cite{herbert2025chronoflux}. It takes one conserved four-dimensional current whose only law is zero divergence and climbs a recovery ladder to inertia, relativity, gravity, the quantum, and cosmology, with a consistency matrix that demands one parameter assignment fit every arena at once. Vibe builds that same current from the discrete side, so the two run the same program from opposite ends. The relation is firm at the bottom rung, and above it the program is a target. Chronoflux is single-author, unreviewed, and unreplicated, so we treat its structure and its discipline as serious and its specific predictions as unconfirmed. It is one program among several, not the subject of the paper.

The gravity-from-thermodynamics programs supply the next family. Einstein's equation comes out as a thermodynamic equation of state from horizon entropy and vacuum entanglement equilibrium \cite{jacobson1995, jacobson2016}, and as the counting of surface degrees of freedom with the cosmological constant left as an integration constant \cite{padmanabhan2010}. Both derivations need a finite count of states in a bounded region, which vibe's discreteness supplies for free with no tuning. These are firm fits. Entropic gravity and its emergent dark sector are the closest published cousin of vibe's wake \cite{verlinde2011}, but the derivation is contested, so it is inspiration and parallel, not a settled rung, and its warning that the law cannot depend on mass alone is a real constraint vibe must respect.

The entanglement-geometry program states that spacetime connectivity is entanglement, with geometry pinching off as the entanglement between regions is cut \cite{vanraamsdonk2010}. Vibe reproduces the co-variation on the substrate, where two regions entangle only through a geometric mediator and the mutual information dies when the channel is severed. The relation is firm on the shape and parallel on the route, since the substrate result uses free-fermion proxies on a graph rather than entanglement generated by the ternary rule itself.

The quantum reconstruction programs set the checklist the emergent quantum surface must pass. Quantum theory is recovered as the unique structure meeting operational axioms of purification and information processing \cite{chiribella2011}, or from a short list of reasonable axioms on outcomes \cite{hardy2001, masanes2011}, with the toy-model line marking the boundary between quantum and classical statistics \cite{spekkens2007}. Vibe reaches the Born rule and the Tsirelson bound by a discrete, deterministic route, which is a parallel demonstration, not a derivation of the continuum reconstruction. The superdeterminism defense is vibe's own stance, framed statistically \cite{hossenfelder2020}. The natively non-classical correlation shape is kept from the work of Christian while its local-realism claim is rejected \cite{christian2007}.

The causal-set and discrete-gravity programs build or constrain emergent spacetime. The causal-set program reaches exact Lorentz invariance only by random sprinkling of points \cite{bombelli1987, dowker2004}, the very move vibe forbids, which makes Lorentz on a fixed lattice vibe's sharpest open wall. Causal dynamical triangulations give a discrete-to-continuum route to four-dimensional spacetime with dimensional reduction at small scales \cite{ambjorn2004}, relational and graph-based quantum gravity supply the relational stance \cite{rovelli1996, konopka2008}, and causal invariance reappears as general covariance in the rewriting-system models \cite{gorard2020}. These are parallels with one shared open tension at the lattice.

The discrete-rule programs put a local reversible rule directly under physics. The cellular-automaton interpretation derives quantum behavior from a deterministic rule and names the two hard walls of Lorentz and locality \cite{thooft2016}. The Dirac equation is recovered from a quantum cellular automaton with a doubly-special route to Lorentz \cite{dariano2014}, and the same equation falls out of discrete sign-alternation and discrete walks \cite{kauffman1996, ord2002, nottale1992}. Reversible conservative logic and the lattice-gas tradition are the direct ancestors of vibe's rule, and four-dimensional isotropy is exactly what forces the 24-cell coin \cite{fredkin1982, frisch1986, dhumieres1986, kari2005}. The hyperbolic-honeycomb substrate, its cell-addressing, and the idea of computing on it draw on Margenstern's two volumes \cite{margenstern2007vol1, margenstern2008vol2}. These are firm program-level kin, with the same open walls.

Two further families bracket the climb. The exceptional-algebra programs reach $D_4$, the 24-cell, the octonions, and the gauge group from the top \cite{baez2002octonions, gunaydin1973, furey2018}, on the algebraic ceiling of the four normed division algebras \cite{hurwitz1898} and the geometry of the 24-cell and its symmetry \cite{coxeter1973, conway2003}. Lisi's E8 proposal is the most ambitious of these, fitting every known particle into the single exceptional group $E_8$ \cite{lisi2007}. Vibe builds the same nucleus from the bottom. The shared open wall is triality to three chiral generations, bound by the Distler-Garibaldi no-go, which proved that no faithful three-generation structure embeds in $E_8$ the way the proposal needs \cite{distler2010}, a wall that also limits the grand-unified and left-right routes \cite{georgi1974, boyle2020}. The upper-layer programs supply the tools for self and scale. Causal emergence gives the test and the lossy-middle constraint \cite{hoel2013}, the Markov-blanket and observer-boundary work gives the self \cite{friston2010, fields2018}, discrete scale relativity gives the self-similar tower \cite{oldershaw2007}, and constructor theory supplies the method and the gravity-induced-entanglement witness \cite{deutsch2015, marletto2017, bose2017}.

This recurring pattern is the paper's underlying argument. Independent workers from string holography, division algebras, cellular automata, causal sets, and quantum foundations keep arriving at the same small set of objects. One conserved local primitive, the exceptional algebra, geometry from entanglement, and a finite state count. Vibe supplies the one thing each of them assumes, a concrete discrete deterministic base from which these objects are derived rather than posited.
